\documentclass{article}
\usepackage{amsmath,hyperref}
\title{Containing a second fault while a controller switch is outstanding}
\begin{document}
\maketitle

\paragraph{Result.}
Adding another client to the same controller-manager switch service
does not give it an independent emergency path. In the pinned manager
the first service callback holds the service-wide mutex while it waits.
The native update loop already provides another path: after an
eligible controller reports an update error, its configured fallback
can be selected before the old service call returns. An independent
hardware read error can also trigger deactivation when the pending
switch skips the failed controller's own update. Neither path,
just by returning an active or inactive software state, certifies
the physical mode. This pins a stronger baseline for I1 rather than
claiming a new fallback mechanism.

\paragraph{Source and experimental boundary.}
The pinned \href{https://github.com/ros-controls/ros2_control/blob/0c94b008a82731967ba89e2e7379ec6d2f07580b/controller_manager/src/controller_manager.cpp}{ros2\_control
manager} uses a shared \texttt{services\_lock\_} in the switch callback,
held for all of \texttt{switch\_controller\_cb}. In non-ASAP switching
that operation waits for the update loop's ready flag before invoking
the switch. The update loop's error handling precedes publishing that
ready flag. Our C++ harness executes the exact service wrapper,
non-ASAP skip predicate, update-error collection and fallback tail,
the resource-overlap and interface-list helpers, the hardware-mode
change helper and the hardware-read branch. A condition variable
holds the old call at the modeled wait, with the native service mutex
still held. The second caller is started while the old
one is outstanding. The inner switch and the hardware and controller
lifecycle effects are stubs. The harness neither exercises the
entire ROS executor nor times any physical actuator.

\paragraph{Same cut, five variations.}
Initially A and B are driving, while C holds a separate interface.
A has an older non-ASAP switch outstanding. By the first update
cut only A's service call has entered its switch body, though the
second client asks to disable B. B's native update error collection
enables the following cases:
\begin{enumerate}
\item With separate scopes, the loop deactivates B's driving controller
and activates its configured hold controller. C remains active.
The recorded mode request was accepted by the stub.
\item If the hold controller declares C's interface as well, the
fallback path also deactivates C. The new dependency is the
fallback's interface, even though C had no error and was not
in either requested service operation.
\item If the hardware stub refuses the mode change, the native helper
logs the refusal and returns no success predicate to the fallback
tail. That path still calls lifecycle deactivation and activation.
The same selected lifecycle bits can therefore accompany accepted
or refused mode transitions. This does not assert how any particular
driver or actual activation subsequently responds.
\item If A's pending switch includes B in the old deactivation
scope, the native non-ASAP guard skips B's update. That cut does
not observe the injected B update error or invoke its fallback.
This is a statement about the skipped cut, not a claim that
B cannot be deactivated by the older switch later.
\item With that same skip, a separate read-side error reporting B's
hardware component leads to B's deactivation before readiness.
That path follows the cached controller-to-hardware mapping, does
not install B's hold controller, and leaves C's controller active
in this fixture. The correct mapping and hardware report are
supplied inputs, not guaranteed by the trace.
\end{enumerate}
The captured state is before the modeled old switch is released.
The standard administrative operations can act after release;
the study does not address whether an old prepared request could
subsequently interfere with a local mode change. In particular,
deactivation at the cut is not a proof that old granted authority
is permanently retired. An update or hardware-read result can
also be late or unavailable; neither fast path is credited in
that branch.

\paragraph{Scope law and comparator.}
For the update path, under the supplied interface resolution,
the selected deactivation set consists of the failed controller
chains together with currently active controllers whose declared
command interfaces intersect any configured fallback interface.
The helper computes this closure before activating the fallbacks.
A healthy active controller outside the union is therefore not
deactivated by that selection. The included Lean facts check
that set criterion, not manager synchronization or device
semantics. Sustained service for C additionally requires a
successful effect path and no latent physical dependency on
the affected interfaces. A list of unaffected controller names
does not establish those facts.

The remote-only timing control is deliberately weaker than a
stock deployment with local fault eligibility, scoped fallback
and independent read evidence. Use the latter as the comparison
where those mechanisms are available. A future end-to-end transfer
would need an independently visible effect certificate and a check
that the old request cannot invalidate it. The selected replay
shows a real alternative to waiting on the shared service, but
does not establish an I1 architecture advance over the existing
local mechanisms or a new physical operating envelope.
\end{document}
